\section{Threats to Validity}\label{sec:threats}

\paragraph{Internal validity.} The cell-permutation test is a global-null permutation: with real corridor effects in the pool it is conservative for the strongest corridors, and it treats cells of one size stratum as equally noisy. Its strictest variant is resolution-limited at $10^{5}$ permutations (Section~\ref{sec:audit}); a ten-fold finer run, which rejects nothing either, is reported in the appendix of the journal version. Cells in different corridors that share a trip stay independent in the scheme, which BY control is meant to cover. The as-of rule prevents the leak we found in a streaming replay (Section~\ref{sec:replay}), but any further replay-based measurement inherits the risk. A leg created just before a rolling-origin window may finish inside it, so its label is not strictly known at the cut.

\paragraph{External validity.} One operator, one country and 22 calendar days of leg creation, in one season of 2018. Whether the corridors flagged persist over months or across carriers is untested; our persistence result is within the 22 days and the same recurring services run in both halves. The data record that a leg ran late, not why.

\paragraph{Construct validity.} The planned time is the dataset's recorded OSRM estimate; we cannot know the version, the profile or whether the carrier's own planning had corridor history, so the error measured is that of this recorded number and not of a production planner. A corridor is an ordered facility pair, which is not a physical road: two services on the same pair may take different roads.

\paragraph{Statistical conclusion validity.} Test legs span six calendar days, so the day-block interval rests on six blocks; the six-hour-block and Diebold--Mariano results agree for the reported model, but six blocks is few. One chronological split and four rolling folds (Section~\ref{sec:rolling}) are the evidence of stability. The validation cut favoured the bar over every Spark model while the test split favoured the models; the reported model was selected on validation by a rule fixed in advance, and both are shown. The 13 slices are not corrected for multiplicity in the figure; the headline claim uses the overall interval. The dependence-aware audit gives zero at its pre-registered strictness: this is a limit of what these data can confirm, not a proof that the corridors are null.

\paragraph{Evaluator and model choice.} The reported model is single-node and the served model is a distributed one; the paper does not claim they are the same. Hyperparameter search was identical in budget across libraries (three learning rates), so a larger search could change the ordering of the close candidates (29.66 to 29.75).
